\chapter[2025 November]{November 2025}

\section[2025/11/27]{Thursday, 27 November 2025}
\label{sec:20251127}

Again bad with checkings not going to do this properly just going to add points form for now

\begin{compactitem}
  \item get nightly build of pytorch to support sm 120 (since I upgraded to 5070 ti) 
  \item update trained to save checkpoints this saves the optimizer state at different epochs so training can be resumed from that point (haven't tested this yet)
  \item saves final model
  \item evaluation script to load model and test it on test data
\end{compactitem}

I'm simplifying a lot but this is the first main test of this

Training 

\renewcommand{\runid}{e4a0fe23-80d7-41b7-8cec-465d24ae9d1c}

\subsection{Run: \texttt{\runid}}

\lstinputlisting[
    caption={Hyperparameters for run \texttt{\runid}},
    label={code:hyperparameters:\runid}
]{figures/\runid/hyperparameters.json}

\begin{figure}[H]
  \centering
  \includegraphics[width=0.48\linewidth]{figures/\runid/embeddings_epoch_25.png}
  \includegraphics[width=0.48\linewidth]{figures/\runid/embeddings_epoch_50.png}
  \includegraphics[width=0.48\linewidth]{figures/\runid/embeddings_epoch_75.png}
  \includegraphics[width=0.48\linewidth]{figures/\runid/embeddings_epoch_100.png}
  \caption{Embeddings at various epochs for run \texttt{\runid}.}
  \label{fig:embeddings:\runid}
\end{figure}

\begin{figure}[H]
  \centering
  \includegraphics[width=0.9\linewidth]{figures/\runid/training_summary.png}
  \caption{Training summary for run \texttt{\runid}.}
  \label{fig:summary:\runid}
\end{figure}

The training looks okay overall nothing special I'm probably going to have to do more hyperparameter tuning, run for longer,
get a more diverse dataset etc but this is good for now to test the pipeline

I created a three\_persons\_test dataset with 3 people

I then ran an evaluation which does the following:

\begin{compactitem}
  \item Loads the saved model checkpoint from training
  \item Runs inference on the test set to generate embeddings for each keypoint
  \item Clusters the embeddings using HDBSCAN (a density-based clustering algorithm that automatically determines the number of clusters)
  \item Visualizes each test image with side-by-side comparisons: predicted clusters (left) vs ground truth (right), with skeletons colored by assignment
  \item Computes clustering quality metrics: Adjusted Rand Index (ARI) and Normalized Mutual Information (NMI)
\end{compactitem}

ARI (Adjusted Rand Index) measures how well the predicted clusters match the ground truth assignments.
It looks at all pairs of points and checks whether they are in the same or different clusters in both the prediction and ground truth.
The "adjusted" part means it accounts for chance - random clustering would score around 0.
The formula is:

\begin{equation}
ARI = \frac{RI - Expected(RI)}{max(RI) - Expected(RI)}
\label{eq:ari}
\end{equation}

where $RI$ (Rand Index) is the proportion of pairs that are either in the same cluster in both assignments or in different clusters in both.
ARI ranges from -1 to 1, where 1 is perfect agreement, 0 is random, and negative means worse than random.

NMI (Normalized Mutual Information) takes an information-theoretic approach.
It measures how much knowing the predicted cluster tells you about the ground truth cluster (and vice versa).
Mutual information between two clusterings $U$ and $V$ is:

\begin{equation}
MI(U, V) = \sum_{i=1}^{|U|} \sum_{j=1}^{|V|} P(i, j) \log \frac{P(i, j)}{P(i)P(j)}
\label{eq:mi}
\end{equation}

where $P(i, j)$ is the probability that a point belongs to cluster $i$ in $U$ and cluster $j$ in $V$.
This gets normalized by the entropy of both clusterings to give a value between 0 and 1:

\begin{equation}
NMI(U, V) = \frac{2 \cdot MI(U, V)}{H(U) + H(V)}
\label{eq:nmi}
\end{equation}

where $H(U)$ and $H(V)$ are the entropies of the two clusterings.
NMI of 1 means the clusters perfectly predict each other, 0 means they share no information.

A more simplified version of this (from claude because this is a bit much) is:

ARI (Adjusted Rand Index) and NMI (Normalized Mutual Information) both measure
how well predicted clusters match ground truth, just from different angles.

ARI looks at pairs of keypoints and asks: if two keypoints belong to the same
person, did they end up in the same cluster? And if they belong to different
people, did they end up in different clusters? It counts how often this is 
right, adjusted so that random guessing scores 0. A score of 1 means perfect
clustering.

NMI asks: if I tell you which cluster a keypoint is in, how much does that tell
you about which person it belongs to? It's based on information theory - a
score of 1 means the clusters perfectly identify the people, 0 means they tell
you nothing.

The evaluation results on 150 test images showed:

\begin{compactitem} 
  \item \textbf{Average clusters found: 3.36} - HDBSCAN found
    an average of 3.36 clusters per image. Since the ground truth is 3 people per
    image, this means it's slightly over-clustering. Some people's keypoints are
    likely getting split into multiple clusters.

  \item \textbf{Average noise points: 1.37} - On average, 1.37 keypoints per
    image couldn't be confidently assigned to any cluster. With 3 people $\times$
    17 keypoints = 51 total keypoints per image, this is only $\sim$2.7\% noise
    which is good.

    \item \textbf{Cluster count accuracy: 51.33\%} - In just over half the images,
      HDBSCAN found exactly 3 clusters matching the ground truth. This is a strict
      metric since getting 2 or 4 clusters counts as wrong even if the overall
      grouping is reasonable.

    \item \textbf{Average ARI: 0.7991} - A score of 0.80 is quite good. It means
      the clustering largely agrees with ground truth even when it doesn't get the
      exact number of clusters right.

    \item \textbf{Average NMI: 0.8434} - A score of 0.84 is strong. The clusters
      are capturing most of the person identity information. 
\end{compactitem}

Overall the embeddings are clearly separating by person as shown by the high ARI and NMI scores.
The main issue is HDBSCAN sometimes over-segments or under-segments, but the underlying embeddings are good.
This could potentially be improved by tuning HDBSCAN's \texttt{min\_cluster\_size} and \texttt{min\_samples} parameters.

Here are some images from the evaluations

\begin{figure}[H]
  \centering
  \includegraphics[width=1\linewidth]{figures/\runid/evaluation/test_image_1.png}
  \includegraphics[width=1\linewidth]{figures/\runid/evaluation/test_image_2.png}
  \includegraphics[width=1\linewidth]{figures/\runid/evaluation/test_image_3.png}
  \caption{Example evaluations for run \texttt{\runid}.}
  \label{fig:evaluations:\runid}
\end{figure}

The first image is an example of it working well, the second is an okay result with some confusion and the third is it working badly

The next experiment was to use the same model but for four people to see how well it generalizes

\begin{compactitem}
\item \textbf{Average clusters found: 4.37} - HDBSCAN found an average of 4.37 clusters per image when the ground truth is 4 people. Still slightly over-clustering but closer to the target than the 3-person case relatively speaking.

\item \textbf{Average noise points: 3.02} - More than double the noise compared to the 3-person test (1.37). With 4 people $\times$ 17 keypoints = 68 total keypoints, this is $\sim$4.4\% noise. The model struggles more with the additional person.

\item \textbf{Cluster count accuracy: 38.67\%} - Dropped from 51.33\% on 3-person data. The model was trained on 3-person scenes so it's not surprising it struggles to generalise to 4 people.

\item \textbf{Average ARI: 0.7969} - Still high despite poor visual results. ARI measures pairwise agreement, so if most keypoints from the same person end up together (even if split across clusters), it still scores well.

\item \textbf{Average NMI: 0.8677} - Actually higher than the 3-person test (0.8434). This is likely because with more people there's more information to capture, and the metric is normalized.
\end{compactitem}

The high ARI/NMI scores are somewhat misleading here. These metrics measure
pairwise relationships - whether keypoints that belong together tend to cluster
together - but they don't heavily penalise over-segmentation or small errors.
Visually the results look worse because even a few misassigned keypoints break
the skeleton coherence. The drop in cluster count accuracy (from 51\% to 39\%)
and the increase in noise points are better indicators that the model doesn't
generalise well from 3-person training data to 4-person scenes. This suggests
training on more diverse data (varying numbers of people) would help.

Here are some images from the evaluations

\begin{figure}[H]
  \centering
  \includegraphics[width=1\linewidth]{figures/\runid/evaluation/test_image_4.png}
  \includegraphics[width=1\linewidth]{figures/\runid/evaluation/test_image_5.png}
  \includegraphics[width=1\linewidth]{figures/\runid/evaluation/test_image_6.png}
  \caption{Example evaluations for run \texttt{\runid}.}
  \label{fig:evaluations:\runid}
\end{figure}

Image 1 is an okay result with some confusion, image 2 and 3 get progressively worse 

The next step is to train on four persons and see how that works and then see how it generalizes to three persons

\renewcommand{\runid}{b1482612-89d8-479d-8faa-4fcc80aabf48}

This is for a "mixed" run where I put 2,3 and 4 persons in the training data

\subsection{Run: \texttt{\runid}}

\lstinputlisting[
    caption={Hyperparameters for run \texttt{\runid}},
    label={code:hyperparameters:\runid}
]{figures/\runid/hyperparameters.json}

\begin{figure}[H]
  \centering
  \includegraphics[width=0.48\linewidth]{figures/\runid/embeddings_epoch_25.png}
  \includegraphics[width=0.48\linewidth]{figures/\runid/embeddings_epoch_50.png}
  \includegraphics[width=0.48\linewidth]{figures/\runid/embeddings_epoch_75.png}
  \includegraphics[width=0.48\linewidth]{figures/\runid/embeddings_epoch_100.png}
  \caption{Embeddings at 25, 50, 75, 100 epochs for run \texttt{\runid}.}
  \label{fig:embeddings_1:\runid}
\end{figure}

\begin{figure}[H]
  \centering
  \includegraphics[width=0.48\linewidth]{figures/\runid/embeddings_epoch_125.png}
  \includegraphics[width=0.48\linewidth]{figures/\runid/embeddings_epoch_150.png}
  \includegraphics[width=0.48\linewidth]{figures/\runid/embeddings_epoch_175.png}
  \includegraphics[width=0.48\linewidth]{figures/\runid/embeddings_epoch_200.png}
  \caption{Embeddings at 125, 150, 175, 200 epochs for run \texttt{\runid}.}
  \label{fig:embeddings_2:\runid}
\end{figure}

\begin{figure}[H]
  \centering
  \includegraphics[width=0.48\linewidth]{figures/\runid/embeddings_epoch_225.png}
  \includegraphics[width=0.48\linewidth]{figures/\runid/embeddings_epoch_250.png}
  \includegraphics[width=0.48\linewidth]{figures/\runid/embeddings_epoch_275.png}
  \includegraphics[width=0.48\linewidth]{figures/\runid/embeddings_epoch_300.png}
  \caption{Embeddings at 225, 250, 275, 300 epochs for run \texttt{\runid}.}
  \label{fig:embeddings_3:\runid}
\end{figure}

\begin{figure}[H]
  \centering
  \includegraphics[width=0.9\linewidth]{figures/\runid/training_summary.png}
  \caption{Training summary for run \texttt{\runid}.}
  \label{fig:summary:\runid}
\end{figure}

I didn't do this perfectly but it is okay for now. I take random images and do the PCA
this was okay when the entire dataset had the same number of people but since it's mixed now
maybe not as applicable but anyways the training looks decent. Interestingly it seems
like at around 100 epochs we have a minimum and then it gets worse. 

While that's the case I'm still going to use the final model (at 300 epochs) for evaluation
now and then maybe try with the 100 and see how much improvement I get

Evaluation on the two persons test set

\begin{compactitem}
\item \textbf{Average clusters found: 2.02} - Almost exactly 2 clusters per image, matching the ground truth perfectly.

\item \textbf{Average noise points: 0.01} - Essentially zero noise. HDBSCAN confidently assigned nearly every keypoint to a cluster.

\item \textbf{Cluster count accuracy: 98.00\%} - The model found exactly 2 clusters in 98\% of images. A massive improvement from the earlier experiments.

\item \textbf{Average ARI: 0.9954} - Near perfect. The predicted clusters almost exactly match the ground truth assignments.

\item \textbf{Average NMI: 0.9963} - Near perfect. The clusters capture virtually all of the person identity information.
\end{compactitem}

Some example images 

\begin{figure}[H]
  \centering
  \includegraphics[width=1\linewidth]{figures/\runid/evaluation/test_image_1.png}
  \includegraphics[width=1\linewidth]{figures/\runid/evaluation/test_image_2.png}
  \caption{Example evaluations for run on the 2 person test set \texttt{\runid}.}
  \label{fig:evaluations_2_persons_test:\runid}
\end{figure}

The first image is perfect (like the majority of them) and the second is an example of some "confusion" when clustering

\todoredefined{I save all the images in the same folder so rerunning evaluation overwrites them - need to fix that later}

Evaluation on the three persons test set

\begin{compactitem}
\item \textbf{Average clusters found: 3.45} - Slightly over-clustering, finding about 3.5 clusters instead of 3. Some people are getting split into multiple clusters.

\item \textbf{Average noise points: 0.53} - Still very low noise, less than 1 keypoint per image on average.

\item \textbf{Cluster count accuracy: 66.00\%} - Improved from the earlier 51.33\% when trained only on 3-person data. The mixed training helped.

\item \textbf{Average ARI: 0.9456} - Strong improvement from 0.7991. The clustering is much more accurate.

\item \textbf{Average NMI: 0.9624} - Up from 0.8434. The clusters capture most of the person identity information.
\end{compactitem}

some example images

\begin{figure}[H]
  \centering
  \includegraphics[width=1\linewidth]{figures/\runid/evaluation/test_image_3.png}
  \includegraphics[width=1\linewidth]{figures/\runid/evaluation/test_image_4.png}
  \caption{Example evaluations for run on the 3 person test set (image 1 and 2) \texttt{\runid}.}
  \label{fig:evaluations_3_persons_test_1:\runid}
\end{figure}

\begin{figure}[H]
  \centering
  \includegraphics[width=1\linewidth]{figures/\runid/evaluation/test_image_5.png}
  \includegraphics[width=1\linewidth]{figures/\runid/evaluation/test_image_6.png}
  \caption{Example evaluations for run on the 3 person test set (image 3 and 4) \texttt{\runid}.}
  \label{fig:evaluations_3_persons_test_2:\runid}
\end{figure}

Image 1 is an example of it performing as expected, image 2 is an example of it performing extremely well
image 3 is an example of it performing poorly but the red blue clusters can be explained the pink, light blue 
and then noise on the figure on the left is a problem and finally the 4th image is quite bad  

Onto the 4 person test set:

\begin{compactitem}
\item \textbf{Average clusters found: 4.45} - Similar pattern to 3-person, slightly over-clustering by about 0.5 clusters on average.

\item \textbf{Average noise points: 0.38} - Very low noise, even with more people in the scene.

\item \textbf{Cluster count accuracy: 65.33\%} - Similar to 3-person (66\%). The model handles both equally well.

\item \textbf{Average ARI: 0.9656} - Excellent. Much better than the 0.7969 when trained only on 3-person data and tested on 4-person.

\item \textbf{Average NMI: 0.9772} - Near perfect information capture despite the increased scene complexity.
\end{compactitem}


\begin{figure}[H]
  \centering
  \includegraphics[width=1\linewidth]{figures/\runid/evaluation/test_image_7.png}
  \includegraphics[width=1\linewidth]{figures/\runid/evaluation/test_image_8.png}
  \caption{Example evaluations for run on the 4 person test set (image 1 and 2) \texttt{\runid}.}
  \label{fig:evaluations_4_persons_test_1:\runid}
\end{figure}

\begin{figure}[H]
  \centering
  \includegraphics[width=1\linewidth]{figures/\runid/evaluation/test_image_9.png}
  \includegraphics[width=1\linewidth]{figures/\runid/evaluation/test_image_10.png}
  \caption{Example evaluations for run on the 4 person test set (image 3 and 4) \texttt{\runid}.}
  \label{fig:evaluations_4_persons_test_2:\runid}
\end{figure}

Image 1 is an example of it performing as expected, image 2 is an example of it performing badly (the clustering is a mess),
Image 3 is decent it seems to get the people right but has some noise on the extremities, Image 4 is a fairly impressive result

Overall the mixed dataset training was a success. The model generalises well
across 2, 3, and 4 person scenes with high ARI and NMI scores throughout. The
2-person test set performed nearly perfectly (98\% cluster accuracy), while 3
and 4 person scenes both hover around 65-66\% cluster accuracy but with
excellent ARI/NMI above 0.94.

The consistent over-clustering (finding ~0.5 extra clusters on average)
suggests the issue might be with HDBSCAN's parameters rather than the
embeddings themselves. If the embeddings were poor we'd see low ARI/NMI, but
they're both above 0.94. What's likely happening is that some people's
keypoints are forming slightly separated sub-clusters in embedding space
(perhaps upper body vs lower body, or occluded limbs), and HDBSCAN is splitting
them. Tuning \texttt{min\_cluster\_size} or using a different clustering
approach that expects a specific number of clusters (like the joint-counting
heuristic mentioned earlier) could help close this gap.

To test this I'm going to try run 4 person evaluation using the model at 100 epochs
which seemed to be the best for training

\begin{compactitem}
\item \textbf{Average clusters found: 4.28} - Closer to 4 than the 300 epoch model (4.45). Less over-clustering.

\item \textbf{Average noise points: 0.37} - Similar to the 300 epoch model. No change here.

\item \textbf{Cluster count accuracy: 75.33\%} - A 10\% improvement over the 300 epoch model (65.33\%).

\item \textbf{Average ARI: 0.9794} - Up from 0.9656. Better cluster assignments.

\item \textbf{Average NMI: 0.9857} - Up from 0.9772. More information captured.
\end{compactitem}

The 100 epoch model outperforms the 300 epoch model across all metrics. This
confirms what the training curves suggested - the model reached optimal
separation around epoch 100 and then started to overfit. The embeddings at 300
epochs are slightly worse which leads to more over-clustering. This highlights
the need for early stopping or a validation set to catch the best checkpoint
rather than just training to completion. Maybe it's not just HDBSCAN's fault which 
is what I was hoping for. A combination of both could produce fairly good results